%%% Research Diary - Entry
%%% Template by Mikhail Klassen, April 2013
%%% 
\documentclass[11pt,letterpaper]{article}

\newcommand{\workingDate}{\textsc{2022 $|$ Outubro}}
\newcommand{\userName}{Prof. Júlio F. C. Santos}
\newcommand{\institution}{IERI-UFU}
\usepackage{researchdiary_png}
% To add your univeristy logo to the upper right, simply
% upload a file named "logo.png" using the files menu above.

\begin{document}
\univlogo

\title{Notas de Aula}

{\Huge Título para a Nota de Aula}\\[3mm]

\textit{Essas notas foram criadas com o intuito de auxiliar o estudo da disciplina de Tópicos em Mercados Financeiro, do Bacharelado em Ciências Econômicas da Universidade Federal de Uberlândia. Para qualquer erro ou omissão, e-mail de contato: \href{mailto:julio.costa@ufu.br}{julio.costa@ufu.br}}

\section*{Renda Variável - Ganhos de Diversificação}


\par Essas notas de aula compilam algumas ideias que se originam do trabalho de \cite{markowitz1952}, \cite{Lintner1965}, \cite{markowitz1976markowitz}  e posteriormente culminaram na formulação do modelo CAPM, desenvolvido por \cite{sharpe1964}.

\par Exemplo de equação numerada:
\begin{equation}
E(R_{P}) = \frac{1}{N}\sum_{i=1}^N E(R_i)
\end{equation}




\begin{figure}[!htb]
    \centering
    \caption{Exemplo figura 1}
    \includegraphics[width=0.7\textwidth]{Ufu_logo.svg.png}
    \caption*{Fonte: Autoria própria.} %% isso é levemente uma gambiarra, mas funciona para o propósito desse template.
    \label{fig 2}
\end{figure}


\subsection*{Exemplo Subseção 1}


\par Exemplo parágrafo para destacar itens:
\begin{itemize}
    \item Exemplo item 1.
    \item Exemplo item 2.
\end{itemize}

%%% as referências devem estar em formato bibTeX no arquivo referencias.bib

\subsection*{Exemplo de Subseção 2}




\printbibliography


\end{document}